% !TeX program = xelatex
% Hamilton-Jacobi wavefronts and trajectories for a free isotropic system.
% Compile: xelatex hamilton_jacobi_tikz.tex
\documentclass[tikz,border=6pt]{standalone}
\usepackage{amsmath}
\usepackage[UTF8,fontset=fandol]{ctex}
\usepackage{fix-cm}
\usetikzlibrary{arrows.meta}

\begin{document}
\begin{tikzpicture}[>=Stealth,thick,every node/.style={font=\fontsize{16}{19}\selectfont}]
  \draw[blue!70!black,->] (-.5,0) -- (8.1,0) node[right] {$q_1$};
  \draw[blue!70!black,->] (0,-3.7) -- (0,3.9) node[above] {$q_2$};
  \fill[black] (0,0) circle (2.2pt);

  % S=constant wavefronts.
  \foreach \r in {1.5,2.7,3.9,5.1,6.3} {
    \draw[red!65,domain=-38:38,samples=80,smooth]
      plot ({\r*cos(\x)},{\r*sin(\x)});
  }
  % Trajectories, normal to the wavefronts.
  \foreach \a in {-32,-16,0,16,32} {
    \draw[blue!70!black,->] (0,0) -- (\a:7.0);
  }
  \draw[white,line width=3pt] (0,0) -- (7.0,0);
  \draw[blue!70!black,->] (0,0) -- (7.0,0);

  \coordinate (A) at ({3.9*cos(16)},{3.9*sin(16)});
  \fill[green!45!black] (A) circle (2.2pt);
  \draw[green!45!black,->,line width=1.5pt] (A) -- ++({1.35*cos(16)},{1.35*sin(16)})
    node[above] {$\boldsymbol p=\nabla S$};
  \node[red,anchor=west] at (4.8,2.55) {$S(q_1,q_2,t)=\text{常数}$};
  \node[blue!70!black,anchor=west] at (5.05,-2.75) {经典轨迹};
  \node[black!70,align=left,anchor=west] at (8.55,.25)
    {作用量的等值面如同波前，\\动量指向其法线方向};
\end{tikzpicture}
\end{document}
